\chapter{Распознавание изображений и видео}
\chaptermark{Распознавание}
\label{sec:recognition}

\section{Задача: одно и то же при разных пикселях}

Глава~\ref{sec:sensors} довела изображение до входа машины: около миллиона выходных нейронов глаза, сжатый поток, в котором передаются в основном края и перемены. Но между датчиком и действием есть шаг, о котором мы пока молчали. Кошка на картинке --- это не набор пикселей. Сдвиньте её на полкадра, поверните, отодвиньте, выключите половину света --- почти все пиксели станут другими, а ответ «кошка» должен остаться тем же. Инженер называет это \emph{инвариантностью}: ответ классификатора не должен меняться при сдвиге, масштабе, повороте и освещении.

Трудность хорошо видна на простом опыте. Возьмём одномерную «сетчатку» из $24$ точек и две фигуры по пять точек, которые могут стоять где угодно. Классификатор, который сравнивает вход с образцом, запомненным в одном месте, угадывает в $49$--$52\%$ случаев --- столько же, сколько монетка. Сдвиг ломает сравнение по пикселям полностью. Нужна другая схема.

\section{Конвейер из ступеней}

Как быстро мозг это делает, мы уже знаем из главы~\ref{sec:code}: животное на мгновенно показанной картинке узнаётся примерно за $150\ms$, и сигнал проходит около десятка ступеней, по $10$--$20\ms$ на каждую~\cite{Thorpe1996}. На каждой ступени нейрон успевает выдать ноль или один импульс. Значит, быстрое распознавание --- это один проход по конвейеру, без долгих раздумий и повторов. Вопрос в том, что делает каждая ступень.

Ответ, который подсказали измерения на первых ступенях зрения~\cite{HubelWiesel1962}, состоит из двух чередующихся операций (рис.~\ref{fig:conveyor}).

\emph{Избирательность.} Нейрон ступени $S$ смотрит на небольшой участок входа и срабатывает, только если там есть определённое сочетание частей: вот этот край, а рядом тот. Это И по частям --- пороговый нейрон главы~\ref{sec:glutcomputer}, у которого порог выше, чем даёт любая одна часть.

\emph{Инвариантность.} Нейрон ступени $C$ собирает выходы многих $S$-нейронов, ищущих одно и то же сочетание в разных местах, и срабатывает, если оно нашлось хоть где-то. Это ИЛИ по положениям --- а точнее, выбор наибольшего.

Для одномерной модели обе операции записываются так:
\begin{empheq}[box=\keybox]{equation}
s_{f,p} \;=\; \Bigl[\sum_i t_{f,i}\,x_{p+i} - \theta\Bigr]_+ ,
\qquad
c_f \;=\; \max_p\, s_{f,p} ,
\label{eq:scstage}
\end{empheq}
где $x$ --- вход, $t_{f,i}$ --- образец признака $f$, $p$ --- положение, $\theta$ --- порог. Первое выражение делает ответ избирательным, второе --- независимым от положения. Наибольшее из многих входов сеть получает из средств главы~\ref{sec:gaba}: взаимное торможение оставляет самый сильный вход (утверждение~\ref{prop:wta}), а деление на общую активность выравнивает ответы при разной яркости~\cite{Carandini2012}.

\begin{figure}[t]
\centering
\begin{tikzpicture}[>=Stealth,
  px/.style={draw,minimum size=0.36cm,inner sep=0pt},
  on/.style={px,fill=blue!55!black},
  snode/.style={circle,draw,very thick,minimum size=0.62cm,inner sep=0pt,font=\scriptsize,fill=blue!6},
  cnode/.style={circle,draw,very thick,minimum size=0.8cm,inner sep=0pt,font=\scriptsize,fill=red!10}]
% two inputs: the same shape at two positions
\foreach \row/\sh/\lab in {0/1/{кадр 1},1/7/{кадр 2}}{
  \begin{scope}[yshift=-\row*1.0cm]
  \node[left,font=\scriptsize] at (-0.25,0) {\lab};
  \foreach \i in {0,...,13}{\node[px] at (\i*0.4,0) {};}
  \foreach \d in {0,1,3,4}{\node[on] at ({(\sh+\d)*0.4},0) {};}
  \end{scope}}
% S stage: detectors of the shape at several positions
\foreach \k in {0,...,5}{
  \node[snode] (S\k) at ({0.8+0.8*\k},1.7) {$S$};}
\node[font=\scriptsize,left] at (-0.25,1.7) {ступень $S$: И по частям};
\draw[->,thick,blue!60!black] (1.2,0.25) -- (S1.south);
\draw[->,thick,orange!85!black] (3.6,0.25) -- (S4.south);
\node[font=\scriptsize,blue!60!black,anchor=east] at (1.25,0.85) {кадр 1};
\node[font=\scriptsize,orange!85!black,anchor=west] at (3.9,0.85) {кадр 2};
% C stage
\node[cnode] (C) at (2.8,3.25) {$C$};
\node[font=\scriptsize,left] at (-0.25,3.25) {ступень $C$: ИЛИ по положениям};
\foreach \k in {0,...,5}{\draw[->,gray!60!black] (S\k.north) -- (C);}
\draw[->,very thick,red!70!black] (C.north) -- ++(0,0.6) node[above,font=\scriptsize,black] {«фигура есть» --- в обоих кадрах};
% next stages
\draw[decorate,decoration={brace,amplitude=5pt},gray!60!black] (6.3,1.4) -- (6.3,-1.3);
\node[right,font=\scriptsize,align=left] at (6.5,0.05) {один и тот же\\объект, другие\\пиксели};
\draw[->,thick,densely dashed,gray!60!black] (3.6,3.25) -- (5.9,3.25) node[right,font=\scriptsize,black,align=left] {следующие пары\\$S$, $C$ --- крупнее,\\сложнее, инвариантнее};
\end{tikzpicture}
\caption{Одна пара ступеней конвейера. Нейроны $S$ ищут одно и то же сочетание частей в разных местах; нейрон $C$ отвечает, если оно нашлось хоть где-то~\eqref{eq:scstage}. Фигура в кадре 1 и в кадре 2 возбуждает разные $S$-нейроны, но один и тот же $C$. Следующие пары ступеней повторяют то же на всё более крупных и сложных сочетаниях.}
\label{fig:conveyor}
\end{figure}

В той же одномерной модели пара ступеней~\eqref{eq:scstage} узнаёт фигуру в любом месте без ошибок, а если каждую точку входа случайно перевернуть с вероятностью $0.1$ --- в $86\%$ случаев, с вероятностью $0.2$ --- в $70\%$. Монетка по-прежнему даёт половину. Сложите несколько таких пар одну над другой: первая ищет края, следующая --- сочетания краёв, ещё выше --- части предметов, наверху --- предметы целиком. С каждой парой сочетания крупнее, а ответ всё меньше зависит от того, где и как предмет лежит~\cite{Riesenhuber1999}.

\begin{keyblock}
\begin{proposition}[Конвейер избирательности и инвариантности]
\label{prop:conveyor}
Быстрое распознавание --- один проход по десятку ступеней. На каждой паре ступеней И по частям~\eqref{eq:scstage} делает ответ избирательным, а ИЛИ по положениям --- независимым от сдвига. Чередование этих двух операций даёт ответ, который различает предметы и почти не зависит от того, где и как они видны. Устройство первых ступеней и скорость прохода измерены; то, что вся цепочка сводится к этим двум операциям, --- упрощение.
\end{proposition}
\end{keyblock}

Если это устройство кажется знакомым, так и должно быть. Именно его чередование --- поиск признака во всех местах, потом объединение по соседним местам --- инженеры перенесли в искусственные \emph{свёрточные} сети~\cite{Fukushima1980}. А недавно пошли обратно: свёрточные сети, обученные распознавать предметы, оказались лучшей из известных моделей ответов нейронов на верхних ступенях зрения~\cite{Yamins2014}. Результат наверху читается просто: по суммарным частотам нескольких сотен нейронов последней ступени линейный решающий блок узнаёт предмет через десятую долю секунды после показа~\cite{Hung2005}. Это частотный код группы из главы~\ref{sec:code}.

\section{Учитель, которого нет: учиться на видео}

Свёрточную сеть учат на миллионах картинок с подписями. Мозгу подписей почти никто не даёт: ребёнок видит предметы часами, а слово «чашка» слышит изредка. Откуда тогда ступени $C$ знают, какие $S$-нейроны объединять? Ведь чтобы объединить «эту фигуру здесь» с «этой фигурой там», надо уже знать, что это одна и та же фигура.

Подсказку даёт само видео. Мир меняется непрерывно: предмет в поле зрения остаётся тем же предметом, пока он сдвигается, поворачивается и меняет освещение, и только иногда взгляд переходит на другой. Всё, что меняется \emph{медленно}, скорее всего относится к предмету, а всё, что меняется быстро, --- к его положению и виду~\cite{WiskottSejnowski2002}. Значит, $C$-нейрону достаточно правила Хебба из главы~\ref{sec:dofa} со следом: связывать то, что шло подряд, а не только то, что шло одновременно~\cite{Foldiak1991}:
\begin{empheq}[box=\keybox]{equation}
\tau_{\text{сл}}\,\frac{\dd\bar y_k}{\dd t} \;=\; -\bar y_k + y_k ,
\qquad
\frac{\dd\mathbf w_k}{\dd t} \;=\; \eta\,\bar y_k\,\bigl(\mathbf x - \mathbf w_k\bigr) .
\label{eq:traceinvariance}
\end{empheq}
Здесь $y_k$ --- активность $C$-нейрона $k$, который выиграл соревнование через торможение, $\bar y_k$ --- её след, та же RC-цепь, что и в правиле~\eqref{eq:threefactor}, а $\mathbf x$ --- выходы $S$-нейронов. Нейрон, выигравший на первом виде предмета, ещё помнит это, когда приходит второй вид, и подтягивает к себе и его.

\begin{figure}[t]
\centering
\begin{tikzpicture}[>=Stealth,x=1.0cm,y=1.0cm]
\foreach \i/\lab in {0/1,1/2,2/3,3/4}{
  \node[draw,thick,fill=blue!8,minimum width=0.9cm,minimum height=0.55cm,font=\scriptsize] at (\i+0.5,2.2) {$A$ в \lab};}
\foreach \i/\lab in {0/4,1/3,2/2,3/1}{
  \node[draw,thick,fill=orange!12,minimum width=0.9cm,minimum height=0.55cm,font=\scriptsize] at (\i+6.5,2.2) {$B$ в \lab};}
\node[font=\scriptsize,gray!40!black] at (5.0,2.2) {пауза};
\draw[->,gray!70!black] (0,0) -- (10.4,0) node[below left,font=\scriptsize,black] {время};
% traces
\draw[thick,blue!60!black] (0,0.05) .. controls (0.4,1.2) and (0.8,1.3) .. (1.2,1.35) -- (4,1.4) .. controls (4.6,0.5) and (5.2,0.15) .. (6,0.08) -- (10,0.05);
\draw[thick,orange!85!black] (0,0.05) -- (6,0.05) .. controls (6.4,1.2) and (6.8,1.3) .. (7.2,1.35) -- (10,1.4);
\node[font=\scriptsize,blue!60!black,anchor=west] at (1.4,1.65) {след $\bar y_1$ --- держится весь проход $A$};
\node[font=\scriptsize,orange!85!black,anchor=west] at (7.2,1.65) {след $\bar y_2$};
\draw[decorate,decoration={brace,amplitude=4pt},gray!60!black] (4.0,2.6) -- (0,2.6);
\draw[decorate,decoration={brace,amplitude=4pt,mirror},gray!60!black] (0,-0.35) -- (4,-0.35) node[midway,below=4pt,font=\scriptsize,black] {все виды $A$ связываются с нейроном 1};
\end{tikzpicture}
\caption{Как видео учит инвариантности, схема. Предмет $A$ проходит через четыре положения, потом пауза, потом предмет $B$. След~\eqref{eq:traceinvariance} нейрона, выигравшего на первом виде $A$, держится весь проход, и все виды $A$ оказываются связаны с одним нейроном. Пауза стирает след, и $B$ достаётся другому нейрону.}
\label{fig:tracevideo}
\end{figure}

Мы проверили это на модели (рис.~\ref{fig:tracevideo}). Два $C$-нейрона соревнуются за входы от $16$ $S$-нейронов: две фигуры в восьми положениях. Фигура проезжает через все положения, по $50\ms$ на каждое, потом полсекунды пауза и следующая случайная фигура. Никаких подписей. Если след короткий, $\tau_{\text{сл}}=5\ms$, нейроны делят вход по \emph{положению}, и ответ совпадает с фигурой не лучше монетки: $52\%$ в среднем по десяти прогонам. Если след длиннее одного положения, $\tau_{\text{сл}}=50$--$200\ms$, во всех десяти прогонах каждый нейрон собирает \emph{все} положения своей фигуры: инвариантность выучена на одном только видео. Пауза между предметами важна: без неё след перетекает на следующий предмет и путает их.

То же самое видно в опыте. Если в эксперименте незаметно подменять предмет, пока глаз перескакивает с одного места на другое, то нейроны верхних ступеней за час-другой начинают отвечать на подменённый предмет как на тот же самый: они связывают то, что шло подряд~\cite{LiDiCarlo2008}. Инвариантность в мозге действительно учится на непрерывности времени. Насколько этим правилом объясняется всё обучение зрения, --- гипотеза.

\section{Видео --- это движение}

Видео --- не только поток кадров для обучения, но и сама задача: что движется, куда и как быстро. Первый шаг нам знаком --- детектор направления главы~\ref{sec:gaba} (рис.~\ref{fig:direction}), в котором запаздывающее торможение гасит движение в одну сторону. Такие детекторы стоят по всему полю зрения, и дальше с ними поступают так же, как с признаками формы: ступени, собирающие много детекторов одного направления в разных местах, отвечают на движение крупного предмета независимо от того, где он. Это та же пара И и ИЛИ~\eqref{eq:scstage}, только признак не «край», а «край, сдвинувшийся за время $d$».

Видео ещё и предсказуемо: следующий кадр почти такой же, как предыдущий. По гипотезе \emph{предсказывающего кодирования} верхние ступени посылают нижним предсказание, а вверх уходит главным образом ошибка --- то, чего предсказание не учло~\cite{RaoBallard1999}. Это та же экономия импульсов, что в утверждении~\ref{prop:economy}: платить за новости, а не за то, что и так известно. Отдельные наблюдения с этой гипотезой согласуются, но как общее устройство конвейера она не доказана.

\section{Мозг и свёрточная сеть}

Насколько далеко заходит сходство? Для быстрого распознавания одного предмета оно действительно глубокое~\cite{Yamins2014}. Но у мозга есть и то, чего нет у простой свёрточной сети.

\emph{Обратные связи.} Конвейер мозга не только прямой: у каждой ступени есть связи назад и вбок. Для трудных картинок --- с перекрытием, в плохом свете --- ответ верхних ступеней складывается позже, и эти поздние ответы лучше описываются сетями с обратными связями, чем прямыми~\cite{Kar2019}. Один проход решает лёгкие случаи, трудные требуют нескольких кругов.

\emph{Устойчивость.} Искусственную сеть можно обмануть незаметной для глаза подделкой пикселей: изменение, которое человек не видит, превращает «панду» в «гиббона»~\cite{Szegedy2014}. Зрение человека к таким подделкам почти нечувствительно. Почему --- вопрос открытый; кандидаты --- шум, деление на общую активность и обратные связи.

\emph{Учитель.} Сети нужны миллионы подписанных примеров, мозгу --- в основном видео без подписей~\eqref{eq:traceinvariance} и немного подсказок от \nt{ДОФА} о том, что важно.

\emph{Энергия.} Весь мозг --- около $20$ ватт (утверждение~\ref{prop:economy}), а распознавание занимает лишь часть этого; обученная свёрточная сеть на графическом ускорителе тратит сотни ватт.

\section{Иллюзии: где машина ошибается и почему}
\label{sec:illusions}

Инженер, который хочет понять чужую схему, подаёт на неё необычные сигналы и смотрит, где она ошибается. Для зрения такие сигналы давно придуманы --- это зрительные иллюзии. Иллюзия --- не поломка. Каждая указывает на определённый механизм машины, который в обычной обстановке работает правильно, а на специально подобранной картинке выдаёт себя.

\emph{Разумные ожидания.} Вход шумный, и машина дополняет его тем, что обычно бывает в мире. Медленные движения встречаются чаще быстрых; когда вход ненадёжен --- например, у картинки мало контраста, --- оценка скорости сдвигается к медленной, и бледный предмет кажется движущимся медленнее яркого~\cite{Weiss2002}. Так же объясняют многие иллюзии яркости: мы видим не яркость пикселя, а то, какой поверхности такая яркость чаще всего соответствовала в прошлом~\cite{Purves2011}. Фотография платья, которое одни видят бело-золотым, а другие сине-чёрным, --- тот же механизм: картинка двусмысленна, и люди расходятся в том, каким они неосознанно считают освещение~\cite{LaferSousa2015,Wallisch2017}. Различия между людьми измерены; то, что мозг здесь сочетает вход с ожиданием по правилам теории вероятностей, --- хорошо обоснованная гипотеза.

\emph{Регулировка усиления.} Посмотрите минуту на водопад, потом на неподвижные камни: камни поплывут вверх. Каналы «вниз» и «вверх» --- пара проводов, как в~\eqref{eq:dualrail}; регулировка усиления~\eqref{eq:nakarushton} приглушила уставший канал, и равновесие пары сдвинулось. Иллюзии наклона и контраста, в которых окружение меняет вид центра, объясняются делением на активность соседей~\cite{Schwartz2009}, как в главе~\ref{sec:gaba}. Есть и поучительный пример осторожности: тёмные пятна на перекрёстках белых полос решётки долго объясняли простым торможением соседей, но опыты с изменённой решёткой показали, что это объяснение не работает~\cite{SchillerCarvey2005}.

\emph{Компенсация задержек.} Путь от глаза до верхних ступеней занимает десятки миллисекунд, и движущийся предмет за это время успевает сдвинуться. Если рядом с ним вспыхнет неподвижная точка, предмет кажется впереди вспышки, хотя они совпадали~\cite{Nijhawan1994}. По одному из объяснений машина продлевает движение вперёд, чтобы скомпенсировать задержку, --- то же, что в главе~\ref{sec:muscles}: при запаздывающей обратной связи приходится предсказывать. Объяснений у этой иллюзии несколько, и спор не решён.

\emph{Ошибка в коде времени.} Неподвижная картинка из колец с резкими цветовыми переходами кажется вращающейся. Контрастные участки обрабатываются быстрее бледных, и разница задержек~\eqref{eq:latency} читается детекторами направления как движение~\cite{Conway2005}. Запускают иллюзию мелкие непроизвольные скачки глаза и моргания: каждый скачок обновляет вход, и детекторы снова видят «движение»~\cite{OteroMillan2012}. Это код временем из главы~\ref{sec:code}, прочитанный неправильно.

\emph{Две картинки в одной.} Каркас куба, который то смотрит на нас одной гранью, то другой, или разные картинки в двух глазах: восприятие перескакивает каждые несколько секунд. Лучшие модели --- взаимное торможение двух толкований, медленная усталость и шум~\cite{MorenoBote2007}, то есть та самая схема~\eqref{eq:halfcentre}, которая в главе~\ref{sec:muscles} порождала ритм шага. Время смены задаёт усталость.

А искусственные сети? Сеть, обученная предсказывать следующий кадр видео, видит вращение на тех же неподвижных кольцах и не видит его на контрольных картинках~\cite{Watanabe2018}. Сети, обученные очищать изображения от шума и повышать их резкость, поддаются тем же иллюзиям цвета и яркости, что и люди~\cite{GomezVilla2020}. Значит, часть иллюзий --- следствие самой задачи, которой учится машина, а не особенностей нервной ткани. Какие иллюзии свойственны только мозгу, --- хорошая проверка для любой модели зрения.

\begin{keyblock}
\begin{proposition}[Иллюзии как отпечатки механизмов]
\label{prop:illusions}
Иллюзия возникает, когда механизм, полезный в обычном мире, получает необычный вход: ожидание побеждает ненадёжный вход, регулировка усиления сдвигает пару каналов, компенсация задержки уводит предмет вперёд, разница задержек читается как движение, взаимное торможение с усталостью перескакивает. Каждая иллюзия --- испытательный сигнал, указывающий на свой механизм. Сами иллюзии измерены; их объяснения --- от хорошо обоснованных до спорных.
\end{proposition}
\end{keyblock}

\section{Итог}

\begin{table}[t]
\centering
\footnotesize
\setlength{\tabcolsep}{4pt}
\renewcommand{\arraystretch}{1.15}
\begin{tabular}{>{\raggedright}p{3.0cm}>{\raggedright}p{4.6cm}>{\raggedright\arraybackslash}p{5.2cm}}
\toprule
Что делает машина & Как & Что известно \\
\midrule
узнаёт за $150\ms$ & один проход по десятку ступеней & скорость измерена \\
делает ответ избирательным & И по частям, ступень $S$~\eqref{eq:scstage} & измерено на первых ступенях \\
делает ответ инвариантным & ИЛИ по положениям, ступень $C$; торможение и деление & измерено на первых ступенях; для верхних --- упрощение \\
выдаёт ответ & частоты сотен нейронов последней ступени, линейное чтение & измерено \\
учится без подписей & правило Хебба со следом~\eqref{eq:traceinvariance} на непрерывном видео & подмена предмета меняет ответы --- измерено; как главное правило --- гипотеза \\
видит движение & детекторы направления, затем та же пара И/ИЛИ & измерено \\
экономит на предсказуемом & вверх идёт ошибка предсказания & гипотеза \\
решает трудные случаи & обратные связи, несколько кругов & поздние ответы и роль обратных связей измерены \\
ошибается предсказуемо & иллюзии: ожидания, регулировка усиления, задержки, взаимное торможение с усталостью & иллюзии измерены; объяснения --- от обоснованных до спорных \\
\bottomrule
\end{tabular}
\caption{Как машина распознаёт изображения и видео.}
\label{tab:recognition}
\end{table}

Распознавание собрано из деталей, которые у нас уже были (таблица~\ref{tab:recognition}): порог даёт И по частям, взаимное торможение --- ИЛИ по положениям, деление --- независимость от яркости, правило Хебба со следом --- учителя, которого нет. Инженеры взяли у зрения чередование избирательности и инвариантности и построили на нём свёрточные сети; обратно мозг пока не отдал главного --- умения учиться на видео без подписей почти без энергии.

\section{Задачи}

\begin{exercise}[\lvlA]
\label{ex:12:1}
\emph{Частота на одной ступени.} Ступень конвейера длится $T=15\ms$, нейроны работают с частотой $r=20$ в секунду. В главе~\ref{sec:code} сто нейронов с независимыми шумами давали за такое время точность около двадцати процентов.
\par\noindent(а)~Какова по~\eqref{eq:rateerror} относительная ошибка частоты одного нейрона за время ступени?
\par\noindent(б)~Теперь шумы коррелированы, $c=0.1$~\eqref{eq:correlated}. Какова точность группы из ста нейронов и каков предел при сколь угодно большой группе? Сколько уровней частоты различимо на одной ступени?
\par\noindent(в)~Что отсюда следует о том, как ступени передают друг другу числа, если корреляция похожа на сигнал? Назовите два выхода из главы~\ref{sec:code}.
\end{exercise}

\begin{exercise}[\lvlA]
\label{ex:12:2}
\emph{Энергия одного узнавания.} Пусть в одном проходе участвуют десять ступеней по $10^7$ нейронов, и каждый выдаёт не больше одного импульса ценой $E_{\text{имп}}\sim10^{-10}$~Дж (глава~\ref{sec:code}).
\par\noindent(а)~Сколько энергии стоит одно узнавание? Какую долю это составляет от энергии всего мозга и от энергии его сигналов~\eqref{eq:energy} за те же $150\ms$?
\par\noindent(б)~С какой средней частотой работают эти нейроны во время прохода? Совместимо ли это со средней частотой~\eqref{eq:energy}, и при каком условии?
\par\noindent(в)~Ускоритель мощностью $300$~Вт узнаёт картинку за $10\ms$. Во сколько раз больше энергии он тратит на одно узнавание?
\end{exercise}

\begin{exercise}[\lvlB]
\label{ex:12:3}
\emph{Пороги ступени $S$.} Одномерная «сетчатка» из $N=24$ точек, фигуры $A=11011$ и $B=10101$ из текста. Образец признака в~\eqref{eq:scstage}: $t_{f,i}=+1$, где у фигуры $1$, и $t_{f,i}=-1$, где $0$ (отрицательный вес --- через \nt{ГАМК}-посредника); $n_f$ --- число единиц фигуры.
\par\noindent(а)~Докажите, что при $n_f-1\le\theta_f<n_f$ нейрон $s_{f,p}$ отвечает тогда и только тогда, когда в окне $p$ стоит в точности фигура $f$. Найдите пороги для $A$ и $B$.
\par\noindent(б)~Найдите пороги, при которых нейрон прощает одну перевёрнутую точку в окне. Проверьте, что при них детектор $A$ не отвечает на чистую $B$ ни в каком положении, и наоборот.
\par\noindent(в)~Докажите, что $c_f$ не меняется при сдвиге фигуры, пока она остаётся внутри «сетчатки».
\par\noindent(г)~Сосчитайте нейроны двух ступеней, если отрицательный вес от каждой точки даёт один общий \nt{ГАМК}-посредник этой точки. Сколько $S$-нейронов понадобится для десяти признаков размером $5\times5$ на картинке $100\times100$?
\end{exercise}

\begin{exercise}[\lvlB]
\label{ex:12:4}
\emph{Сколько длится одна картинка из двух.} Перескоки восприятия описывают схемой~\eqref{eq:halfcentre}.
\par\noindent(а)~Пусть $\tau\ll\tau_a$. Пока побеждает группа~1, $r_2=0$ и $r_1=I-a_1$. Найдите $a_1(t)$ и $a_2(t)$ и выведите условие, при котором группа~2 вырывается вперёд.
\par\noindent(б)~В симметричном цикле усталость каждой группы в начале и в конце её победы одна и та же. Составьте уравнения для длительности победы $T_d$ и решите их численно при $I=1$, $\beta=2$, $b=2$, $\tau_a=0.4\,\text{с}$.
\par\noindent(в)~Смоделируйте~\eqref{eq:halfcentre} при $\tau=20\ms$ и при $\tau=2\ms$ и сравните период с (б) и с рис.~\ref{fig:halfcentre}. Откуда разница?
\par\noindent(г)~Докажите, что в пределе (а) $T_d$ пропорционально $\tau_a$ и не зависит от $I$. Какое $\tau_a$ даёт смену каждые три секунды? Проверьте моделированием.
\end{exercise}

\begin{exercise}[\lvlB]
\label{ex:12:5}
\emph{Моделирование: цена инвариантности.} Используйте функцию \texttt{two\_stage} из \texttt{models/\allowbreak recognition\_models.py} ($4000$ попыток, \texttt{seed}$\,=0$).
\par\noindent(а)~При $N=24$ постройте долю верных ответов обоих классификаторов в зависимости от вероятности переворота точки $p$ от $0$ до $0.5$. При каком $p$ пара ступеней $S$, $C$ ошибается в каждом четвёртом случае?
\par\noindent(б)~При $p=0.1$ измерьте долю верных ответов для $N=8$, $12$, $24$, $48$, $96$, $192$. Как она зависит от $N$?
\par\noindent(в)~Объясните зависимость (б): что делает взятие наибольшего по всё большему числу положений, когда вход шумный? Предложите изменение ступени $C$, которое уменьшило бы эту потерю, и объясните, почему в зрении ступени $C$ собирают лишь соседние положения.
\end{exercise}

\begin{exercise}[\lvlB]
\label{ex:12:6}
\emph{Моделирование: окно следа.} Используйте функцию \texttt{trace\_learning} из \texttt{models/\allowbreak recognition\_models.py} (десять прогонов, \texttt{seed}$\,=0,\dots,9$).
\par\noindent(а)~При паузе $0.3\,\text{с}$ постройте среднюю и наихудшую оценку инвариантности в зависимости от $\tau_{\text{сл}}$ от $5\ms$ до $2\,\text{с}$ (логарифмическая шкала). В каком окне $\tau_{\text{сл}}$ все десять прогонов выучивают инвариантность безошибочно?
\par\noindent(б)~Повторите без паузы.
\par\noindent(в)~Объясните верхнюю границу окна (а) через~\eqref{eq:traceinvariance}: что должно случиться со следом за паузу?
\par\noindent(г)~Проверьте, зависит ли нижняя граница от времени показа одного положения (\texttt{dwell}$\,=0.2\,\text{с}$ вместо $50\ms$), от скорости обучения (\texttt{eta} от $0.5$ до $8$) и от шага интегрирования (\texttt{dt}$\,=1\ms$). Что это говорит о её происхождении?
\end{exercise}

\begin{exercise}[\lvlB]
\label{ex:12:7}
\emph{Проектирование: движение в любом месте.} Ряд из $24$ точек с шагом $\Delta x=0.1^\circ$. На каждой паре соседних точек стоит детектор направления главы~\ref{sec:gaba}: \nt{ГЛУТ}-нейрон $Y$ возбуждается правой точкой $B$ и тормозится левой точкой $A$ через \nt{ГАМК}-посредника с задержкой $d$; торможение гасит возбуждение, если они приходят не дальше чем на $\Delta_{\text{т}}=5\ms$ друг от друга. Над детекторами --- ступень $C$~\eqref{eq:scstage}.
\par\noindent(а)~В какую сторону движения отвечает $Y$? Для какой полосы скоростей одна задержка $d$ гасит противоположное направление?
\par\noindent(б)~Нужно, чтобы $C$ отвечал на движение в свою сторону в любом месте ряда и молчал при движении в противоположную сторону со скоростью от $5$ до $20^\circ$ в секунду. Хватит ли одной задержки? Предложите схему и найдите все допустимые задержки.
\par\noindent(в)~Сосчитайте нейроны схемы. Что произойдёт при скорости $40^\circ$ в секунду в противоположную сторону?
\end{exercise}

\begin{exercise}[\lvlC]
\label{ex:12:8}
\emph{Подделка пикселей.} Раздел «Мозг и свёрточная сеть» оставил открытым вопрос, почему зрение почти нечувствительно к незаметным подделкам.
\par\noindent(а)~В модели \texttt{two\_stage} решает разность $\Delta=c_{\text{верн}}-c_{\text{чуж}}$, где $c_f$ --- наибольшее по окнам $k$ число совпадений фигуры $f$ с окном $k$. Докажите, что для смены ответа нужно не меньше $\lceil(\Delta+1)/2\rceil$ перевёрнутых точек, а для ничьей --- $\lceil\Delta/2\rceil$.
\par\noindent(б)~Найдите $\Delta$ для чистых фигур при $N=24$ и проверьте полным перебором, сколько переворотов на самом деле нужно.
\par\noindent(в)~Классификатор «число единиц на входе больше $3.5$ --- это $A$» тоже инвариантен к сдвигу и безошибочен на чистых картинках. Сколько переворотов нужно, чтобы его обмануть? Какова его доля верных ответов при $p=0.1$ по сравнению с $0.86$ у пары ступеней?
\par\noindent(г)~Открытый вопрос. Что даёт устойчивость к подделке: инвариантность или избирательность? Предложите изменение модели, которое увеличит наименьшее число переворотов, не ухудшив ответа на случайный шум, и проверьте его. Какой из кандидатов текста --- шум, деление на общую активность, обратные связи --- можно проверить в этой модели и как?
\end{exercise}
